% V. Методика. Този раздел описва ПЛАНА на експеримента и не съдържа нито едно число,
% което е резултат. Числата, които се появяват тук, са настройки на опита, зададени
% предварително, и параметри на средата.
\section{Методика за експериментална проверка}
\label{sec:method}

\subsection{Какво се проверява}

Проверяват се три твърдения, които проектните решения от Раздел~\ref{sec:alternatives}
приемат за верни. Първо, че двупосочният канал доставя стойност до клиента по-бързо от
периодичното запитване при еднаква честота на обновяване. Второ, че двата формата за обмен
се различават по обем на съобщението и по време за разбор. Трето, че системата запазва
работоспособност при нарастващ брой едновременни клиенти до определена граница, която
експериментът трябва да намери, а не да предположи.

Всяко от трите твърдения е записано тук \term{преди} провеждането на измерването. Целта е
една погрешна прогноза да остане видима и да бъде отчетена като погрешна, вместо да бъде
пренаписана след получаването на резултата.

\subsection{Измервани величини и начин на измерване}

\textbf{Закъснение.} Измерва се като \term{възраст на стойността при пристигане}: от
момента, в който събирачът произведе извадка, до момента, в който клиентът разчете
съобщението. Всяко съобщение носи времеви отпечатък от сървъра в поле \code{ts}, а
клиентът изважда този отпечатък от собствения си часовник. И двете страни се изпълняват на
една машина, така че въпросът за синхронизация на часовниците отпада.

Източникът на време от двете страни е \code{std::chrono::system\_clock}, приведен до
микросекунди. Изборът на стенен, а не на монотонен часовник е наложен от задачата: двата
края са различни обекти и монотонният часовник няма обща начална точка между тях. Всички
продължителности, измервани изцяло вътре в един процес, ползват
\code{std::chrono::steady\_clock}. Разделителната способност на \code{system\_clock} е
измерена на използваната машина, а не приета наум: най-малката наблюдавана ненулева
стъпка между две последователни четения е 100 ns, тоест поне с три порядъка по-малка от
измерваните закъснения. Самото четене на часовника се извършва веднъж на съобщение и
затова не влиза в измерваната величина.

\textbf{Обем на съобщението.} Измерва се като брой байтове по мрежата за един и същ пакет
със стойности: при WebSocket~\cite{rfc6455} байтовете на кадъра, включително заглавната му
част; при събитията, изпращани от сървъра, байтовете на парчето заедно с рамката на
събитието; при периодичното запитване байтовете на заявката и на пълния отговор, заглавните
части включително. Мери се доставеното по мрежата, а не дължината на полезния товар, защото
именно доставеното се плаща.

\textbf{Време за изписване и за разбор.} Измерва се вътре в процеса, около самото
извикване, с монотонния часовник. И двата анализатора изграждат пълно дърво, за да е
сравнението между еднакъв вид работа, а не между изваждане на едно поле и построяване на
документ.

\textbf{Пропускателна способност при нарастващ брой клиенти.} Измерва се като брой
доставени съобщения за фиксиран прозорец при 1, 8, 32 и 128 едновременни връзки по
WebSocket, съпоставен с очаквания брой, който е произведение на броя клиенти и броя
интервали в прозореца.

\subsection{Инструмент за натоварване}

Натоварването и измерването се извършват от \code{pulse-demo}, изграден от същото
хранилище. Той стартира сървъра в отделна нишка в същия процес, свързва се към него по
локалния интерфейс като обикновен клиент и говори същите три протокола. Клиентската страна
не ползва браузър: браузърният клиент е проверен функционално, но не участва в числата,
защото планировчикът на страницата и честотата на опресняване на екрана биха се смесили с
измерваната величина.

Това решение има две последствия, които трябва да бъдат заявени. Първо, измерващият клиент
и сървърът делят един процес, затова показанието за време на процесора в края на опита
включва и работата на измерващия клиент; то се отчита като показание на процеса, а не като
цена на сървъра. Второ, обменът минава през стека на локалния интерфейс, тоест измерените
закъснения не съдържат мрежово време и представляват долна граница за реална мрежа.

Че самият инструмент не е тясното място, се установява по следния начин: при 128 връзки
броят доставени съобщения се сравнява с очаквания. Ако измерващият клиент изоставаше,
доставените съобщения биха били по-малко от очакваните, а закъснението би нараствало с
времето вътре в прозореца. И двете се отчитат в Раздел~\ref{sec:results}.

\subsection{Условия за валидност на измерването}

Условията са фиксирани предварително и не са променяни, след като са видени резултати.

\begin{enumerate}[topsep=0pt,itemsep=0pt,leftmargin=*]
\item Интервалът на публикуване е 100 милисекунди за всички опити.
\item Изписването и разборът се измерват в 5 кръга по 10\,000 повторения. Отчита се
      средното на \term{средния} кръг, а най-ниският и най-високият кръг се отчитат
      като мярка за разсейване. Един кръг сам по себе си се различава от друг с
      множител, близък до две, затова едно число от един кръг не би било възпроизводимо.
\item За всеки транспорт се събират 65 съобщения, от които първите 5 се отхвърлят.
      Отхвърлените носят еднократните разходи: установяване на връзката, първата извадка
      и страниците, които разпределителят на памет още не е докоснал.
\item Опитът с нарастващ брой клиенти ползва прозорец от 3 секунди за всяка стойност на
      броя клиенти.
\item Отчитат се средна стойност, медиана и деветдесет и пети процентил, а не само
      средна стойност.
\item По време на опита на машината не се изпълнява друго забележимо натоварване, а
      браузърът е затворен.
\end{enumerate}

Целият опит се възпроизвежда с една команда, \code{pulse-demo --bench}, а изведеният от
нея отчет е приложен без редакция в Приложение~Г.

\subsection{Апаратна и програмна среда}

Средата на измерването е част от резултата и е описана изцяло, за да може опитът да бъде
повторен \tabref{tab:environment}.

\begin{table}[H]
\caption{Апаратна и програмна среда на измерването}
\label{tab:environment}
\centering
\fontsize{11}{13.2}\selectfont
\begin{tabular}{@{}L{4.6cm}L{8.4cm}@{}}
\toprule
\textbf{Елемент} & \textbf{Стойност} \\
\midrule
Процесор & AMD Ryzen 5 3600, 6 ядра, 12 логически процесора \\
Оперативна памет & 15,9 GiB \\
Операционна система & Windows 11 Pro N, версия 10.0.26200 \\
Компилатор & g++ 15.2.0, MinGW-w64, x86\_64-posix-seh \\
Стандарт на езика & C++20, без разширения на компилатора \\
Ниво на оптимизация & \code{CMAKE\_BUILD\_TYPE=Release} \\
Система за изграждане & CMake 4.3.2, Ninja 1.13.2 \\
Мрежов път & локален интерфейс, \code{127.0.0.1} \\
Интервал на публикуване & 100 ms \\
Браузър & не участва в измерването, вж. подраздел 5.3 \\
\bottomrule
\end{tabular}
\end{table}
